\begin{frame}{낙관적 초기화: 시작값 하나로 탐험을 유도}
  모든 $Q_0(a)$ 를 실제 가능한 보상보다 \textbf{훨씬 크게}(낙관적으로)
  초기화.
  \vskip0.5em
  \begin{itemize}
    \item 아직 안 당겨본 팔은 ``과대평가된'' 상태 $\to$
      순수 \textbf{탐욕적 선택만}으로도 자연스럽게 먼저 시도.
    \item 시도해서 실망하면(실제 보상 $<$ 초기값) 추정치가 내려가고,
      그다음 팔로 넘어감.
    \item \textbf{초반 몇 번의 강제적 탐험을 낙관적 시작값 하나로}
      만들어내는 셈.
  \end{itemize}
  \vskip0.4em
  \begin{block}{왜 ``크게'' 잡아야 하는가}
    작은 값(예: 0)으로 초기화하면 증분 갱신식이 추정치를
    \textbf{올리는 방향}으로만 움직임 $\to$ 어떤 팔이든
    ``0보다 낫네''로 판정받은 순간 나머지 팔이 \textbf{영원히
    외면}될 수 있음. \textbf{과대평가(상단에서 시작)}만이
    ``기회는 실망하면 회수한다''는 올바른 방향의 탐험을 만듦.
  \end{block}
\end{frame}
